% TOP_PROBLEM: {"id": 20001033, "title": "Exact trace certificates and a regular counterexample normal form for Caccetta--Haggkvist", "category_id": 2, "category": {"id": 2, "name": "combinatorics", "display_name": "Combinatorics", "description": "Counting problems, graph theory, discrete structures.", "slug": "combinatorics", "order_index": 2, "created_at": "2026-07-31T15:26:25.671Z"}, "status": "partially_solved", "problem_number": "AIM-COMBINATORICS-0158", "set_id": 14}
% FIRST_POSTED: 2026-10-01
% ENTRY_KIND: statement_only
% UnsolvedMath ID 20001033; source catalogue code AIM-COMBINATORICS-0158
% !TeX program = lualatex
% Definition and sources only; this file is not an LLM solution attempt.
\documentclass[11pt,a4paper]{article}
\usepackage[margin=25mm]{geometry}
\usepackage{fontspec}
\setmainfont{lmroman10-regular.otf}[BoldFont=lmroman10-bold.otf,
 ItalicFont=lmroman10-italic.otf,BoldItalicFont=lmroman10-bolditalic.otf]
\usepackage{amsmath,amssymb,mathtools}
\usepackage{unicode-math}
\setmathfont{latinmodern-math.otf}
\AtBeginDocument{\let\setminus\smallsetminus}
\usepackage{xurl,hyperref}
\hypersetup{colorlinks=true,urlcolor=blue}
\setcounter{secnumdepth}{0}
\setlength{\parindent}{0pt}
\setlength{\parskip}{5pt}
\setlength{\emergencystretch}{3em}
\begin{document}
\section*{Exact trace certificates and a regular counterexample normal form for Caccetta--Haggkvist}\label{20001033}
{\small UnsolvedMath ID 20001033 · AIM-COMBINATORICS-0158 · Combinatorics · AIM Workshop Problem Lists}
\subsection{Definitions and mathematical statement}
Let $D=(V,E)$ be a finite simple loopless directed graph on $n$ vertices.
There is at most one arc $u\to v$ for an ordered pair of distinct
vertices; the two opposite arcs between a pair are allowed.
The outdegree $d^+(v)$ counts arcs leaving $v$, and
$\delta^+(D)=\min_{v\in V}d^+(v)$.
A directed cycle of length $\ell$ is a sequence of distinct vertices
$v_0,\ldots,v_{\ell-1}$ with arcs $v_i\to v_{i+1}$, indices modulo
$\ell$. In particular, opposite arcs form a cycle of length two.

The Caccetta--H\"aggkvist conjecture asserts that for every integer
$r\ge1$, every such nonempty digraph with $\delta^+(D)\ge r$ has a
directed cycle of length at most $\lceil n/r\rceil$.
Equivalently, if its shortest directed cycle has length $g$, then
$n\ge r(g-1)+1$. Positive minimum outdegree ensures that some cycle
exists, by following successive outgoing arcs in a finite graph.

The source also gives a matrix formulation for the digon-free case.
Let $A$ be a zero--one $n\times n$ matrix with diagonal entries one
and $a_{ij}a_{ji}=0$ for $i\ne j$. If every row sum is at least
$r+1$, the proposed conclusion is
$\operatorname{tr}(A^{\lceil n/r\rceil})>n$.
The diagonal ones allow a closed walk to be padded to that length;
the surplus trace detects a nonconstant closed walk and hence a short
cycle in the underlying loopless digraph.
\subsection{Short English statement}
Does minimum outdegree at least r force a directed cycle of length at most the ceiling of n divided by r?
\subsection{Sources}
\begin{itemize}
\item[S1] ULAM.AI and the UnsolvedMath Contributors, supplied problems.json, record 20001033 (AIM-COMBINATORICS-0158), AIM Workshop Problem Lists. Catalogue identity and source transcription; CC BY 4.0.
\url{https://huggingface.co/datasets/ulamai/UnsolvedMath}.
\item[S2] American Institute of Mathematics, The Caccetta-Haggkvist conjecture, item 2. Original workshop question underlying AIM-COMBINATORICS-0158.
\url{https://aimath.org/WWN/caccetta/caccetta.pdf}.
\end{itemize}
\end{document}
